\chapter{Kinds of Obstacle and Their Dependencies}
\label{ch:05}

The preceding chapters described the coupling among crises and the disciplines in which they are studied. This chapter sorts the problems by the kind of obstacle each presents and by their dependence on one another. The program begins from the problems that human communities already know they need to address, and not from the assumption that a new class of system will arrive to address them. This ordering has a methodological purpose. When a technology is introduced first and its applications are sought afterward, the technology tends to become an explanatory placeholder, a term that stands in for the unspecified mechanism by which difficulties are supposed to be resolved. The interview remarks on this in connection with claims that advanced systems will end poverty and halt climate change, observing that the logic by which a language model accomplishes the latter is never supplied. An inventory-first method forces each intended application to be stated as a problem with a definite structure, so that one can ask what kind of obstacle it presents and what kind of intervention could remove it. The relevant domains are numerous and interdependent: climate instability, biodiversity loss, freshwater scarcity, soil degradation, energy poverty, persistent pollution, infectious disease, senescence, institutional fragility, armed conflict, inequality, educational exclusion, and limitations in the pace and reliability of scientific discovery itself. The term polycrisis is appropriate where the emphasis falls on the interaction of such crises, so that their combined effects exceed what their separate analysis would predict, whereas metacrisis is better reserved for the deeper institutional, epistemic, and civilizational conditions that generate and perpetuate the crises, including the failure of coordination and the degradation of shared methods for establishing what is true.

Within this inventory the first discrimination to be made is among kinds of obstacle. Some problems are scientific, in that what is lacking is knowledge: a catalyst for efficient nitrogen fixation, a mechanism for degrading a persistent pollutant, a materials system for storing energy at the required density, or a design for shielding against radiation at acceptable mass. Others are engineering problems, in which the relevant science is adequate but construction at the required scale, cost, or reliability has not been achieved. A third kind is institutional, concerning the design of rules, incentives, and organizations through which resources are allocated and conflicts are settled. A fourth kind concerns values, where parties disagree about what ought to be done and no further fact will settle the disagreement. A crucial observation is that the categories are not exclusive and that most significant problems are composite. The negotiation of competing claims to a river basin has a hydrological and engineering component that analysis can reduce, an institutional component concerning the forum and procedure of negotiation, and a value component concerning the relative weight of agricultural, urban, ecological, and cultural uses, and this last component is not diminished by any improvement in forecasting. Greater intelligence may help with all of these, by enlarging the set of options, clarifying consequences, and identifying agreements that were not previously visible, but it does not make political disagreement or conflicting values disappear, and a program that treated every problem as scientific would reproduce the very error of misattributed agency that motivates the argument against the dominant narratives.

The second discrimination concerns dependence among problems. The inventory is not a list but a structure, because the resolution of one problem often presupposes the resolution of another, and some problems are enabling in that their solution relaxes constraints on many others. Abundant low-cost energy, for example, changes the feasibility of desalination, of materials recovery, and of large-scale water transport simultaneously, while reliable institutions for verification change the feasibility of every proposal that depends on trusting claims one cannot check directly. The identification of such enabling problems supplies a principled order of effort, and it is also the place at which the distinction among scientific, engineering, institutional, and value components becomes useful, since only the first two respond directly to capability. The program accordingly assigns to capable systems the tasks in which they can supply what is lacking, which are chiefly the scientific and engineering components, and it assigns to human institutions the components in which what is lacking is a decision.

\section*{Formalization}

Let $P$ be a finite set of problems and let $\tau:P\to 2^{\{S,E,I,V\}}$ assign to each problem the nonempty set of component types it contains, where $S$, $E$, $I$, and $V$ denote scientific, engineering, institutional, and value components respectively. Let $G=(P,\to)$ be a directed acyclic graph in which $p\to q$ means that the resolution of $p$ is a prerequisite for the resolution of $q$.

\begin{definition}[Capability-resolvable problem]
A problem $p$ is \emph{capability-resolvable} if $\tau(p)\subseteq\{S,E\}$. A problem is \emph{decision-bearing} if $\tau(p)\cap\{I,V\}\neq\emptyset$.
\end{definition}

For the value component, let a finite set of persons $\{1,\dots,n\}$ hold complete preorders $\succeq_i$ (weak preference relations) over a set $X$ of outcomes. Say that $x$ \emph{Pareto-dominates} $y$ if $x\succeq_iy$ for all $i$ and $x\succ_jy$ for some $j$, and let $\mathrm{PO}(X)$ denote the outcomes in $X$ that no element of $X$ Pareto-dominates. Let $X_c\subseteq X$ denote the outcomes attainable at capability level $c$, with $X_c\subseteq X_{c'}$ for $c\le c'$.

\begin{proposition}
(i) If $x,y$ are distinct elements of $\mathrm{PO}(X_c)$, then either every person is indifferent between $x$ and $y$, or at least one person strictly prefers $x$ to $y$ and at least one person strictly prefers $y$ to $x$. Consequently, whenever $\mathrm{PO}(X_c)$ contains two outcomes between which some person is not indifferent, Pareto efficiency underdetermines collective choice from $X_c$. (ii) Increased capability need not remove the underdetermination: there exist profiles and levels $c<c'$ with $|\mathrm{PO}(X_c)|=|\mathrm{PO}(X_{c'})|=2$.
\end{proposition}

\begin{proof}
(i) Suppose some person strictly prefers $x$ to $y$ and no person strictly prefers $y$ to $x$. Because the preorders are complete, the absence of a strict preference for $y$ means that every person weakly prefers $x$ to $y$, so $x$ Pareto-dominates $y$, contradicting $y\in\mathrm{PO}(X_c)$. The symmetric case is excluded in the same way, so either strict preferences run in both directions or there are none. (ii) Take $n=2$, $X_c=\{x,y\}$ with $x\succ_1y$ and $y\succ_2x$, and let $X_{c'}=X_c\cup\{z\}$ where both persons strictly prefer each of $x$ and $y$ to $z$. Then $z$ is dominated and $\mathrm{PO}(X_{c'})=\{x,y\}$.
\end{proof}

\begin{remark}
The proposition does not assert that every selection rule must originate outside the preferences being aggregated. A social-choice rule such as majority rule, a maximin criterion, or a bargaining solution is a function of the profile. What Pareto efficiency leaves open is the choice among such rules and a procedure for adopting one, and with three or more outcomes and at least two persons Arrow's theorem~\cite{arrow1951} shows that no rule over strict rankings with unrestricted domain satisfies Pareto efficiency, independence of irrelevant alternatives, and non-dictatorship together (Arrow, 1951). That choice is a decision for the legitimate procedure $\Sigma$ and not an output of capability.
\end{remark}

In the dependency graph let $\mathrm{depth}(p)$ denote the length of the longest path ending at $p$ and let $\mathrm{out}(p)$ denote the number of problems reachable from $p$. The \emph{enabling score} of $p$ is $\mathrm{out}(p)$, and the program prioritizes capability-resolvable problems of maximal enabling score, subject to the verification and authority constraints of the chapters that follow.
